% =============================================================================
%  ES827 - Robótica Industrial
%  Resumo: da Cinemática Inversa ao Planejamento de Trajetória
%
%  COMO COMPILAR (no terminal):
%      pdflatex resumo_cinematica_trajetoria.tex
%  (rode 2x para o sumário/índice ficarem corretos). Gera o PDF de mesmo nome.
% =============================================================================
\documentclass[11pt,a4paper]{article}

% ----- Codificação e idioma -------------------------------------------------
\usepackage[utf8]{inputenc}      % permite acentos direto no arquivo (á, ç, ã)
\usepackage[T1]{fontenc}         % fontes que suportam acentos no PDF
\usepackage[brazil]{babel}       % hifenização e termos em português

% ----- Matemática -----------------------------------------------------------
\usepackage{amsmath,amssymb}     % ambientes e símbolos matemáticos
\usepackage{bm}                  % negrito em matemática (vetores/matrizes)

% ----- Layout e cores -------------------------------------------------------
\usepackage[margin=2.3cm]{geometry}
\usepackage{xcolor}
\usepackage{tcolorbox}           % caixas coloridas para destaques
\usepackage{enumitem}
\usepackage{hyperref}

% Cores personalizadas
\definecolor{azul}{HTML}{1F6FB2}
\definecolor{verde}{HTML}{2E8B57}
\definecolor{laranja}{HTML}{C0560E}
\definecolor{vinho}{HTML}{8B1A1A}

\hypersetup{colorlinks=true, linkcolor=azul, urlcolor=azul}

% Comando-atalho: rótulo de aula em cinza
\newcommand{\aula}[1]{{\small\textcolor{gray}{(#1)}}}

\title{\textbf{ES827 — Robótica Industrial}\\[2pt]
       \large Da Cinemática Inversa ao Planejamento de Trajetória}
\author{Resumo de estudo}
\date{\today}

\begin{document}
\maketitle

\begin{tcolorbox}[colback=gray!8,colframe=gray!40,title=Onde você parou]
Logo após a última aula de cinemática inversa \aula{aula 9 — cinemática inversa
iterativa}. Abaixo: primeiro um \textbf{resumo leve} do caminho até ali, depois
a explicação da \textbf{parte seguinte} — planejamento de trajetória \aula{aula 10}.
\end{tcolorbox}

% =============================================================================
\section*{\textcolor{laranja}{Parte 1 — Resumo do que veio antes (aulas 3 a 9)}}
% =============================================================================

\subsection*{1) Transformações homogêneas \aula{aula 3}}
Toda pose (posição + orientação) de um corpo é uma matriz $4\times4$:
\[
T \;=\;
\begin{bmatrix} \bm{R}_{3\times3} & \bm{p} \\[2pt] 0\;\;0\;\;0 & 1 \end{bmatrix},
\qquad
\begin{aligned}
\bm{R} &= \text{rotação }(3\times3)\\
\bm{p} &= \text{posição }(3\times1)
\end{aligned}
\]
Compor movimentos = multiplicar matrizes: $\;T^{0}_{2} = T^{0}_{1}\,T^{1}_{2}$.

\subsection*{2) Denavit--Hartenberg e graus de liberdade \aula{aula 4}}
Cada junta $i$ é descrita por 4 parâmetros DH: $a_i,\ \alpha_i,\ d_i,\ \theta_i$
(os três primeiros são geometria fixa; $\theta_i$ \emph{ou} $d_i$ é a variável da junta).
\[
T^{i-1}_{i} = \mathrm{Rot}(z,\theta_i)\,\mathrm{Trans}(z,d_i)\,
              \mathrm{Trans}(x,a_i)\,\mathrm{Rot}(x,\alpha_i)
\]

\subsection*{3) Cinemática \textsc{direta} \aula{aulas 5--6}}
\emph{Dadas} as juntas $q=(\theta_1,\theta_2,\dots)$, onde está a ferramenta?
\[
T^{0}_{n}(q) = T^{0}_{1}\,T^{1}_{2}\cdots T^{n-1}_{n}
\]
Sempre única e fácil: basta substituir os valores das juntas.

\subsection*{4) Jacobiano e velocidades \aula{aula 7}}
Relaciona a velocidade das juntas $\dot q$ com a do efetuador $\dot X$:
\[
\boxed{\;\dot X = J\,\dot q\;}\qquad
J_{ij} = \frac{\partial X_i}{\partial q_j}
\]
\textbf{Singularidade}: quando $\det(J)=0$, $J$ não tem inversa — o robô
``trava'' uma direção, ou exigiria velocidade/torque infinitos
(tipos: punho, cotovelo, ombro).

\subsection*{5) Cinemática \textsc{inversa} analítica \aula{aula 8}}
Pergunta inversa: dada a pose desejada $X$, quais juntas $q$ a produzem?
\[
X\ (\text{desejado}) \;\longrightarrow\; q = ?
\qquad\Bigl(\text{resolver } T^{0}_{n}(q)=X\Bigr)
\]
É o problema \textbf{difícil}: pode ter várias soluções (cotovelo p/ cima ou
p/ baixo), nenhuma, ou infinitas. Resolve-se por geometria/álgebra.

\subsection*{6) Cinemática \textsc{inversa iterativa} \aula{aula 9} \textmd{\textcolor{vinho}{$\leftarrow$ última aula antes de você se perder}}}
Quando não dá para isolar $q$ na mão, resolve-se \emph{numericamente}
(método de Newton aplicado à robótica):
\begin{enumerate}[leftmargin=2.2em]
  \item Chute inicial $q^{(0)}$.
  \item Erro de pose: $\ \Delta X = X_{\text{des}} - f\!\left(q^{(k)}\right)$
        \quad($f=$ cinemática direta).
  \item Correção: $\ \Delta q = J^{-1}\!\left(q^{(k)}\right)\Delta X$.
  \item Atualiza: $\ q^{(k+1)} = q^{(k)} + \Delta q$.
  \item Repete até $\lVert \Delta X\rVert < \text{tolerância}$.
\end{enumerate}
Ideia: ``ande'' das juntas na direção que reduz o erro, usando $J^{-1}$ como
bússola. Por isso singularidades ($\det J = 0$) quebram o método.

% =============================================================================
\section*{\textcolor{laranja}{Parte 2 — A parte seguinte: Planejamento de Trajetória \aula{aula 10}}}
% =============================================================================

\begin{tcolorbox}[colback=verde!6,colframe=verde!50]
A cinemática inversa responde \emph{``quais juntas levam a um ponto''}.\\
O planejamento de trajetória responde:
\textbf{``como o robô se move, no tempo, de um ponto até outro?''}\\[3pt]
Queremos uma função suave $q(t)$ que dê a posição de cada junta a cada
instante, com velocidade e aceleração bem comportadas.
\end{tcolorbox}

\subsection*{A) Onde isso entra}
$\text{Planejador de Tarefa} \to \boxed{\text{Gerador de Trajetória}} \to
 \text{Controlador} \to \text{Robô}$.
O gerador recebe os pontos-objetivo $X^{*}(t)$ e produz a referência $q^{*}(t)$.

\subsection*{B) Dois espaços para planejar}
\begin{itemize}[leftmargin=1.4em]
  \item \textcolor{azul}{\textbf{Espaço Cartesiano}}: planeja o caminho do
        efetuador no mundo $(x,y,z)$. Previsível (ex.: reta), mas exige IK a
        cada ponto e pode cruzar singularidades.
  \item \textcolor{vinho}{\textbf{Espaço das Juntas}}: planeja direto cada
        $q_i(t)$. Simples e seguro (sem IK em tempo real); caminho cartesiano
        menos previsível. \textbf{É o foco da aula.}
\end{itemize}

\subsection*{C) A ponte com a aula 9 — exemplo do SCARA}
Usa-se a IK iterativa para achar as juntas de \emph{cada} ponto. Cinemática
direta do SCARA:
\[
\begin{Bmatrix} X\\ Y\\ Z \end{Bmatrix}
=
\begin{Bmatrix}
a_2\cos(\theta_1+\theta_2) + a_1\cos\theta_1\\
a_2\operatorname{sen}(\theta_1+\theta_2) + a_1\operatorname{sen}\theta_1\\
d_1 - d_3
\end{Bmatrix}
\]
Constantes: $a_1,a_2,d_1$. \quad Variáveis: $\theta_1,\theta_2,d_3$.
Calcula-se $J$ e $J^{-1}$ e itera-se ($\Delta q = J^{-1}\Delta X$).
Cada pose-objetivo vira um vetor de juntas:

\begin{center}
\begin{tabular}{l|cccc}
& Origem & Posição 1 & Posição 2 & Origem\\\hline
$\theta_1\,[^\circ]$ & 90 & 76{,}95 & 78{,}28 & 90\\
$\theta_2\,[^\circ]$ & $-90$ & $-56{,}25$ & $-110{,}92$ & $-90$\\
$d_3\,[\text{m}]$ & 0{,}1 & 0{,}3 & 0{,}5 & 0{,}1\\
\end{tabular}
\end{center}

\subsection*{D) Atribuir tempos aos movimentos}
$0\text{s} \xrightarrow{1\text{s}} \text{Pos1} \xrightarrow{2\text{s}}
 \text{Pos2} \xrightarrow{2\text{s}} \text{Origem}$ \;(em $t=0,1,3,5$\,s).
Cada junta vira pontos $(t, q)$ a serem ligados por uma \emph{curva suave}:
ligar por retas dá velocidade descontínua (trancos), então usamos polinômios.

% -----------------------------------------------------------------------------
\section*{\textcolor{laranja}{Trajetórias suaves — as 3 ferramentas}}
% -----------------------------------------------------------------------------

\subsection*{1) Polinômio \textsc{cúbico} (3ª ordem) — posição e velocidade}
Tenho $(t_0,q_0)$ e $(t_f,q_f)$ e fixo as velocidades $v_0,v_f$ (em geral 0):
\[
q(t)=a_0+a_1 t+a_2 t^2+a_3 t^3, \qquad
\dot q(t)=a_1+2a_2 t+3a_3 t^2.
\]
4 incógnitas $\Leftrightarrow$ 4 condições:
\[
\begin{bmatrix}
1 & t_0 & t_0^2 & t_0^3\\
0 & 1 & 2t_0 & 3t_0^2\\
1 & t_f & t_f^2 & t_f^3\\
0 & 1 & 2t_f & 3t_f^2
\end{bmatrix}
\begin{bmatrix} a_0\\a_1\\a_2\\a_3 \end{bmatrix}
=
\begin{bmatrix} q_0\\ v_0\\ q_f\\ v_f \end{bmatrix}
\]
\textit{Exemplo} ($\theta_1:10^\circ\!\to\!45^\circ$ em 1\,s, $v=0$):
$\ \theta_1(t)=10+105t^2-70t^3$.

\subsection*{2) Polinômio de \textsc{5ª ordem} (quíntico) — inclui aceleração}
Fixa posição, velocidade \emph{e} aceleração nas duas pontas (6 condições):
\[
\begin{aligned}
q(t)  &= a_0+a_1 t+a_2 t^2+a_3 t^3+a_4 t^4+a_5 t^5\\
\dot q(t) &= a_1+2a_2 t+3a_3 t^2+4a_4 t^3+5a_5 t^4\\
\ddot q(t) &= 2a_2+6a_3 t+12a_4 t^2+20a_5 t^3
\end{aligned}
\]
\textit{Exemplo} ($\theta_2:10^\circ\!\to\!45^\circ$ em 1\,s, $v=a=0$):
$\ \theta_2(t)=10+350t^3-525t^4+210t^5$.

\subsection*{3) LSPB — segmento linear com concordância parabólica}
Perfil de velocidade \textbf{trapezoidal}: acelera, cruzeiro a velocidade
constante, desacelera. 3 trechos:
\[
q(t)=
\begin{cases}
q_0+\dfrac{\alpha}{2}\,t^2, & 0\le t\le t_b \quad(\text{parábola: rampa de }v)\\[8pt]
\dfrac{q_0-q_f+V t_f}{2}+V t, & t_b\le t\le t_f-t_b \quad(\text{reta: }v\text{ const.})\\[8pt]
q_f-\dfrac{\alpha t_f^2}{2}-\dfrac{\alpha}{2}\,t^2, & t_f-t_b< t\le t_f \quad(\text{parábola: freia})
\end{cases}
\]
$\alpha=$ aceleração das pontas, \ $V=$ velocidade de cruzeiro, \ $t_b=$ tempo
de mistura. Posição em ``S'', velocidade em trapézio, aceleração em degraus
$(+\alpha,\,0,\,-\alpha)$.

% -----------------------------------------------------------------------------
\section*{\textcolor{laranja}{Resumão — qual usar?}}
% -----------------------------------------------------------------------------
\begin{tcolorbox}[colback=azul!5,colframe=azul!40]
\begin{tabular}{ll}
\textcolor{azul}{\textbf{Cúbico (3ª)}}   & fixa posição e velocidade. Mais simples.\\
\textcolor{vinho}{\textbf{Quíntico (5ª)}} & fixa também a aceleração $\Rightarrow$ mais suave.\\
\textcolor{verde}{\textbf{LSPB}}          & velocidade constante no meio (trapézio).\\
\end{tabular}
\end{tcolorbox}

\noindent\textbf{Fluxo completo da aula 10:}
pontos no espaço $\to$ \aula{IK iterativa, aula 9} $\to$ juntas $q$ de cada
ponto $\to$ atribui tempos $\to$ \aula{cúbico / quíntico / LSPB} $\to$
$q(t)$ suave. \checkmark

\vspace{4pt}
{\small\textcolor{gray}{Próximas aulas: incertezas (11), dinâmica (12--13),
controle linear (14).}}

\end{document}
